\title{Large Language Models Can Automatically Engineer Features for Few-Shot Tabular Learning}

\begin{abstract}
Large Language Models (LLMs) have demonstrated impressive competence in solving difficult and novel reasoning tasks, positioning them as promising tools for tabular learning—a critical component in many practical scenarios. In this work, we introduce a new in-context learning paradigm, termed \textsf{FeatLLM}, where LLMs are leveraged as feature generators, transforming raw input data into a format optimized for predictive modeling on tabular datasets. The engineered features serve as inputs to a basic machine learning predictor, such as linear regression, which enables robust few-shot learning. At inference, the \textsf{FeatLLM} approach relies solely on these pre-constructed features and the downstream model, avoiding repeated LLM queries. Unlike many current LLM-based solutions, \textsf{FeatLLM} does not necessitate querying the LLM for every test instance during prediction. It also operates effectively with standard API-level access to LLMs, thus circumventing the restrictions imposed by limited prompt lengths. Empirical validation across a diverse array of tabular datasets reveals that \textsf{FeatLLM} systematically produces superior feature transformations, delivering a substantial (average 10\%) improvement in performance over existing methods such as TabLLM and STUNT.
\end{abstract}

\section{Introduction}
Large Language Models (LLMs) have recently distinguished themselves through their proficiency in tackling previously unseen tasks, often delivering strong performance without explicit fine-tuning for each new domain~\cite{creswell2022selection,imani2023mathprompter,taylor2022galactica}. These models, built upon vast and diverse textual corpora, encapsulate extensive knowledge that can often substitute for specialized expertise in a broad range of contexts~\cite{Genomes2021,singhal2023large,wilcox2003role}. Consequently, LLMs have become instrumental in facilitating automation across various fields, encompassing challenges such as dialogue understanding~\cite{finch2023leveraging} and the rectification of software errors~\cite{fan2023automated}. LLMs are particularly effective when provided with descriptive prompts and illustrative examples: this allows them to capture task intent and recognize salient features within data~\cite{brown2020language,zhao2023survey}, thereby functioning as adept feature selectors.

Efforts to capitalize on LLMs’ embedded knowledge and advanced reasoning have expanded into tabular data analysis. As an example, an LLM’s implicit understanding—such as associations between advancing age and susceptibility to certain health conditions—can directly inform medical risk predictions. Contemporary approaches formulate both tasks and features in natural language and transform data instances into text, which are then processed by an LLM~\cite{dinh2022lift,wang2023anypredict}. Subsequent strategies might leverage in-context learning~\cite{nam2023semi} or efficient fine-tuning paradigms~\cite{dinh2022lift,hegselmann2023tabllm} to enhance performance. Empirically, LLMs’ prior knowledge has often led to superior outcomes in low-data scenarios, surpassing traditional tabular modeling baselines~\cite{hegselmann2023tabllm}.

Existing LLM-based methods for learning from tabular data exhibit several notable drawbacks. Achieving end-to-end prediction with these approaches typically demands at least one inference pass per instance, which leads to substantial computational overhead. Furthermore, attaining high predictive accuracy often hinges on fine-tuning the underlying LLMs—a process that is impractical for those models lacking accessible training resources or APIs for tuning. This issue is particularly significant since the latest high-performing LLMs typically restrict user interaction to inference APIs only~\cite{anil2023palm,openai2023gpt4}. Another core problem is that most current strategies do not efficiently accommodate long prompts~\cite{liu2023lost}. When tabular datasets include a large number of features, the text-based representation can exceed feasible lengths, resulting in either technical infeasibility or marked reductions in predictive accuracy. Collectively, these limitations undermine the applicability of LLM-based tabular learning systems in real-world scenarios.

In our work, we depart from relying on LLMs solely for end-to-end decision-making. Rather, we reposition their role to focus on feature construction, thus tapping into their general knowledge in a more strategic fashion. We present a new in-context learning framework termed \textsf{FeatLLM}, which is specifically designed to uncover the “criteria” LLMs implicitly use to form predictions. As an illustration, for a disease prediction task, an LLM is prompted to distill and produce explicit rules that describe which combinations of feature values are indicative of the presence of the disease. These derived criteria or rules are then leveraged to engineer novel features, ultimately enabling replacement of original features in downstream tasks. This paradigm reassigns the primary workload of the LLM to feature generation, greatly simplifying subsequent modeling—once these features are available, inference can be carried out efficiently using lightweight models, thereby reducing inference time. Leveraging LLMs’ in-context learning capabilities, our method enables feature extraction simply by including a series of illustrative examples in the prompt, entirely circumventing the need for conventional model training. Additionally, we utilize principles from ensemble learning—specifically, feature bagging~\cite{breiman2001random}—to address issues associated with lengthy prompts when working with feature-rich datasets.

\textsf{FeatLLM} decomposes the prompt-driven process for generating novel features into two main stages: (1) comprehending the task and inferring how features relate to targets, and (2) employing insights from the first stage along with a small set of demonstrations to formulate concrete rules for class separation. This logical structuring facilitates the LLMs in discarding features that lack discriminatory power during the rule formulation process. The resulting rules are translated to programmatic code, which is then applied to the original data; as a result, each sample is converted into a binary feature encoding whether the rule is satisfied. Subsequently, a lightweight model is trained to map these new binary features to class probability estimates. Model stability is further enhanced using an ensemble technique, whereby the process of feature creation is repeated under feature bagging. By judiciously selecting the number of features and instances in each bag, the risk associated with excessively large prompts is reduced. Notably, \textsf{FeatLLM} does not require gradient-based training and enjoys low cost during inference, making the use of LLMs in a wide range of practical scenarios both feasible and efficient.

We benchmark \textsf{FeatLLM} across 13 distinct tabular datasets under limited data scenarios, demonstrating both robustness and high effectiveness. Our findings indicate that LLMs are capable of leveraging external prior knowledge as well as knowledge learned from the given data to synthesize valuable features. The proposed method consistently surpasses leading few-shot learning approaches across numerous experimental conditions. The code implementation is made available anonymously at~\texttt{https://github.com/Sungwon-Han/FeatLLM}.

\section{Related Work}

\subsection{Few-Shot Learning with Tabular Data} Recent progress in few-shot learning has enabled models to learn representations that are transferable and robust, even when only limited annotations are available~\cite{chen2018closer}. Although early research concentrated primarily on the image domain, the versatility of few-shot techniques has since been demonstrated on a wide array of data types, such as text, audio, and, more recently, tabular datasets~\cite{majumder2022few,nam2023stunt,schick2021exploiting}. The reliance on a very limited number of labeled examples raises concerns regarding the capture of unintended or spurious associations~\cite{bartlett2021deep}. As a result, current work in this field often incorporates auxiliary data sources or embeds domain-specific priors to guide the model's inductive biases during the training process. Notably, semi-supervised methods harness large volumes of available unlabeled data combined with thoughtful data augmentations~\cite{ucar2021subtab,yoon2020vime}, and some approaches advocate for generic model pre-training using unsupervised objectives before adaptation to new tasks~\cite{somepalli2021saint}. Unsupervised training paradigms can entail the intentional masking or corruption of inputs with a reconstruction objective~\cite{arik2021tabnet,majmundar2022met}, or leverage contrastive learning objectives based on various data augmentations~\cite{bahri2022scarf,chen2020simple}. Nevertheless, due to the diverse structure and attributes of tabular data, universal augmentation methods remain elusive, limiting the practical success of these strategies within few-shot regimes~\cite{nam2023stunt}.

More recently, research has explored pre-training models across a multitude of tabular datasets extending far beyond the intended target task, subsequently utilizing transfer learning for downstream adaptation~\cite{zhang2023generative,levin2022transfer,zhu2023xtab}. An example is TransTab~\cite{wang2022transtab}, which leverages transformer models to capture semantic linkages between different columns through the use of column names in addition to raw table entries. Such cross-dataset knowledge transfer supports the deployment of models for new tabular prediction problems in zero-shot and few-shot scenarios. Another noteworthy contribution, TabPFN~\cite{hollmann2022tabpfn}, adopts the strategy of generating synthetic data with assorted distributions to pre-train a Bayesian neural network. Post pre-training, inference for novel tasks involves direct application on their training and test data without fine-tuning. The diversity of learned priors from multiple datasets enabled these models to surpass established tree-based benchmarks, demonstrating tangible benefits for few-shot tabular learning.

\subsection{Language-Interfaced Tabular Learning} An additional avenue for incorporating prior knowledge is found in large language models (LLMs)~\cite{anil2023palm,openai2023gpt4}. Trained on vast and varied textual datasets, LLMs have proven capable of effective transfer and generalization to entirely new tasks, validating their utility across an extensive array of application domains~\cite{brown2020language}. There is a growing body of literature investigating how to adapt the generalized linguistic knowledge embedded in LLMs for problems involving tabular data. Most methodologies involve transforming tables into a serialized textual format, embedding feature and task descriptions within prompts that are then ingested by the LLM~\cite{dinh2022lift,hegselmann2023tabllm,wang2023anypredict}. This technique facilitates the model's ability to identify and leverage relationships among features and tasks solely via prompt construction. Advancements in this area encompass both parameter-efficient finetuning methods such as LoRA~\cite{hu2021lora} and IA3~\cite{liu2022few}, as well as in-context learning approaches, where few-shot demonstrations are supplied to the LLM at inference time in the absence of further training~\cite{dinh2022lift,hegselmann2023tabllm,nam2023semi,slack2023tablet}.

While the previously discussed approaches have shown success in enhancing few-shot tabular learning, their dependency on routing all inference through the LLM introduces difficulties for real-world industrial implementation. To address this limitation, this work shifts away from the standard black-box methodology where every prediction passes end-to-end through the LLM. Instead, our strategy involves instructing the LLM to explicitly determine the distinguishing rules or conditions for each target class. \textsf{FeatLLM} then converts these extracted rules into executable code for feature engineering, thereby facilitating predictions without further reliance on the LLM. This method capitalizes on in-context learning, using prior expertise along with few-shot examples to extract governing rules.

\section{Methods}
The architecture of \textsf{FeatLLM} focuses on distilling the defining “rules” that correspond to each class, utilizing a handful of provided samples as illustrated in Figure~\ref{fig:main_model}. The rules generated by the LLM are programmatically converted into binary feature functions; for any given input, these functions yield indicators reflecting rule satisfaction. These binary features are then aggregated via a dot product with a set of learnable non-negative weights, producing class likelihood estimates. Model training enables the system to quantify the relevance of each feature with respect to class membership (further details are provided in Section~\ref{Sec:training}). This methodology employs bagging to repeatedly generate and utilize ensembles of rules, subsequently integrating their outputs. The ensuing sections detail the overall problem formulation and the nuances of each procedural step.

\vspace*{-0.12in}
\paragraph{Problem formulation.} Consider a tabular dataset containing $N$ labeled instances, denoted as $\mathcal{D} = \{(\mathbf{x}^i, \mathbf{y}^i)\}_{i=1}^{N}$. Each entry in $\mathcal{D}$ consists of $d$ attributes, where $\mathbf{x}^i$ represents a vector of dimensionality $d$, and is accompanied by descriptive natural language feature names, $F = \{f_j\}_{j=1}^{d}$, along with their separate definitions. Within a classification framework, every label $\mathbf{y}^i$ is taken from a finite collection of possible classes. For $k$-shot learning experiments, a randomly chosen subset of $k$ labeled examples, where $k < N$, is utilized for model training.

\subsection{Prompt Design for Extracting Rules}
\label{Sec:prompt_design}
To promote the extraction of rules by the LLM following more reliable reasoning procedures, we structure the prompt to reflect the strategy an expert human would apply in analyzing tabular data. The constructed prompt consists of three principal segments (illustrated in Fig.~\ref{fig:prompt_for_rule} and detailed in Appendix~\ref{sec:example_rule}):

\vspace*{-0.12in}
\paragraph{Basic information description.} This section introduces the vital contextual details needed for problem resolution (see the highlighted orange text in Fig.~\ref{fig:prompt_for_rule}). This encompasses the task description specifying the labels, outlines and definitions of the features, as well as representative example cases. The task itself is framed as a question (e.g., “Does this patient’s myocardial infarction complications data indicate chronic heart failure? Yes or no?”; further illustrations appear in Appendix~\ref{sec:task_desc}). Descriptions for each feature specify data type and may optionally add a concise explanation. Categorical fields can list example category values. For demonstrations, a few labeled instances from the training data are converted into textual form, each coupled with the corresponding ground-truth label. The serialization scheme mirrors approaches adopted in earlier studies~\cite{dinh2022lift,hegselmann2023tabllm}:

\begin{align}
&\text{Serialize}(\mathbf{x}^i,\mathbf{y}^i, F) = \nonumber \\ 
&``\ f_1\ \text{is}\ \mathbf{x}^i_1.\ ... \ f_d\  \text{is}\  \mathbf{x}^i_d.\ \ \text{Answer:}\  \mathbf{y}^i\ ”
\end{align}

\vspace*{-0.12in}
\paragraph{Reasoning instruction.} Our objective is to bolster the reasoning ability of the LLM by embedding reasoning instructions into the prompts (highlighted in blue in Fig.~\ref{fig:prompt_for_rule}). The reasoning guidance starts with a preliminary sentence that mirrors the chain-of-thought technique~\cite{wei2022chain}. Following this, the LLM's inference workflow is segmented into two distinct phases. Initially, the LLM is tasked with establishing the causal links or typical associations between the relevant features and the task's objective, entirely independent of provided examples; this approach leverages the model's foundational knowledge or general reasoning, prompting a more comprehensive extraction of its pre-existing understanding of the problem. Subsequently, in the next phase, the model integrates the earlier causal analysis with specific instance demonstrations to synthesize a set of classification rules for each class. This dual-phase process is designed to reduce the risk of the LLM drawing upon superficial or irrelevant feature correlations, thereby guiding its attention to more pivotal characteristics. To maintain a balance and circumvent extracting either an insufficient or overwhelming quantity of rules—which may cause underfitting or unwieldy prompt lengths—we predetermine a constant number of rules (specifically, 10)\footnote{Analysis on the number of rules can be found in Section~\ref{sec:analysis}.} to be identified per class. Moreover, to ensure rule compatibility, the model is prompted to consider the feature types as detailed in the task's foundational information, thus minimizing the possibility of type-based inconsistencies in the extracted rules.

\vspace*{-0.12in}
\paragraph{Response instruction.} To make the generated rules more amenable to subsequent analysis and implementation, we explicitly instruct the LLM on structuring its answers (see yellow text in Fig.~\ref{fig:prompt_for_rule}). The prompt lays out clear specifications for both the organizational scheme of responses according to answer class (i.e., required response formatting) and the prescribed template for individual rules (i.e., rule formatting). These directives facilitate the LLM's ability to appropriately adapt the rule structure based on each feature's data type. If not otherwise constrained, the LLM is capable of composing more complex rules by logically combining several using connectors like AND or OR. For an illustration of the resulting output, refer to Appendix~\ref{sec:outcome-rule}.

\subsection{Inferring Class Likelihood via Rules}
\label{Sec:training}

\paragraph{Parsing rules for feature generation.} The rules produced in the earlier step serve as the basis for constructing new binary features. For each class, these features reflect whether a sample meets the respective class’s rule set, taking a value of either ‘0’ or ‘1’. Such features can then be utilized for estimating class membership probabilities. Nevertheless, because the rules elicited from the LLM are delivered in natural language, it becomes necessary to engineer a specific parsing mechanism for automated feature extraction. Although we enforce a standardized structure in the formulation of rule-based responses to facilitate parsing, LLM outputs may still contain minor syntactic inconsistencies~\cite{kovavc2023large}. For example, the LLM may occasionally translate operators such as “$>$” or “$<$” to verbal equivalents like “greater than” or “less than,” complicating the parsing process.

To overcome the difficulties presented by noisy or inconsistent text, we choose not to develop sophisticated bespoke parsers; instead, we rely on the LLM to handle the variability. The LLM’s proficiency in semantic comprehension allows it to interpret rules despite stylistic alterations, and its extensive training on code makes it capable of generating concise code upon request. To capitalize on these capabilities, the prompt given to the LLM is structured to include the function name, clear input and output specifications, and the relevant extracted rules. The generated prompt, along with example results, can be found in Figures~\ref{fig:prompt_for_parsing} and ~\ref{sec:example_parsing}-\ref{sec:example_parse_outcome} in the Appendix. Upon receiving the output, the code produced by the LLM is executed directly via Python's \texttt{exec()} to convert rule descriptions into usable binary features.

\vspace*{-0.12in}
\paragraph{Inferring class likelihood.} When binary features representing the satisfaction of specific rules for each class are available, a straightforward way to estimate the likelihood that a sample belongs to a class is by tallying the number of rules satisfied within each class—that is, taking the sum over that class's binary features. Nevertheless, this approach neglects the fact that individual rules and features may contribute unequally; some may be much more indicative of class membership than others. To capture and appropriately weight these contributions, we use a standard machine learning (ML) model to learn the relative importance of features associated with each class. Specifically, we can utilize a linear model without an intercept term. Define the vector of binary features produced from sample $\mathbf{x}^i$ for class $k$ as $\mathbf{z}_k^i \in \{0, 1\}^R$, where $R$ indicates the number of rules per class and $c$ denotes the number of classes. The probability assigned to each class $\mathbf{p}^i$ can then be expressed as:

\begin{align}
\text{logit}_k^i &= \max(\mathbf{w}_k, 0) \cdot \mathbf{z}_k^i, \nonumber \\
\mathbf{p}^i &= \text{Softmax}({[\text{logit}_{1}^i, ..., \text{logit}_{c}^i]}). \label{eq:infer_prob}
\end{align}

In these equations, the vectors $\mathbf{w}_k$ denote trainable parameters that indicate the relevance of each feature within the linear classifier. By enforcing the weights to be non-negative, we guarantee that features extracted by the LLM cannot decrease a class's probability estimate during training. This setup feeds the computed logit values into a softmax layer, translating them to class membership probabilities. Training is conducted using a cross-entropy loss, employing a limited set of labeled samples to adjust the weights $\mathbf{w}_k$.

\vspace*{-0.12in}
\paragraph{Ensembling with bagging.} To further enhance reliability, we repeat the full procedure multiple times, constructing an ensemble of inference models. For $T$ iterations, we collect the learned sets of weights $\{\mathbf{w}[t]\}_{t=1}^T$. To determine the ultimate prediction, we average the probabilities $\mathbf{p}^i$ for each class, each computed using the respective weights $\mathbf{w}[t]$ and following Eq.~\ref{eq:infer_prob}.

To infuse stochasticity into the rule-generation process for the ensemble, we configure LLM inference with a relatively elevated temperature. Additionally, we shuffle the sequence of few-shot examples, as prior work has noted that the ordering of demonstrations can influence the outcomes produced by the LLM~\cite{lu2022fantastically}\footnote{For further discussion on rule diversity, refer to Appendix~\ref{sec:diversity}}. To further harness the diversity inherent in the predictions for greater reliability, we utilize bagging by randomly sampling subsets of either features or instances during each trial. This approach is particularly beneficial as the size of the training set or the number of features increases. The ensemble strategy we propose yields two principal benefits. Firstly, even if the LLM occasionally outputs erroneous rules, the correct rules identified in other ensemble trials can offset these mistakes, thus reinforcing the robustness of the system. This phenomenon relates to the notion of self-consistency in LLMs~\cite{wang2022self}: rules that consistently appear across different trials have a higher likelihood of being valid. Secondly, our ensemble technique mitigates constraints imposed by the limited prompt size of LLMs. With large tabular datasets, when the data surpasses what can be input in a single prompt, bagging serves to alleviate this restriction, helping to preserve the model’s reliable functioning. While any single rule set may fail to encapsulate all feature interactions, the aggregate of multiple weak learners developed across numerous trials can effectively counteract the limitations due to incomplete feature or instance representation in individual cases.

\section{Experiments}
We assess \textsf{FeatLLM} across a variety of tabular datasets and conduct detailed examinations to better understand our framework’s behavior. For brevity, supplementary experiments—including those involving different LLMs, qualitative case studies, and additional evaluations—are provided in the Appendix.

\subsection{Performance Evaluation}
\paragraph{Datasets.}
Our empirical study is based on 13 benchmark datasets for tasks involving binary or multi-class classification: (1) Adult~\cite{asuncion2007uci} (determining if an individual’s annual income exceeds \$50,000); (2) Bank~\cite{moro2014data} (classifying whether a client is likely to subscribe to a term deposit); (3) Blood~\cite{yeh2009knowledge} (predicting donor return behavior for subsequent donations); (4) Car~\cite{kadra2021well} (evaluating the classification of car quality); (5) Communities~\cite{misc_communities_and_crime_183} (forecasting crime rates in various regions); (6) Credit-g~\cite{kadra2021well} (assessing whether an individual represents a high or low credit risk); (7) Diabetes\footnote{\texttt{kaggle.com/datasets/uciml/pima-indians-diabetes-database}} (identifying diabetes presence in patients); (8) Heart\footnote{\texttt{kaggle.com/datasets/fedesoriano/heart-failure-prediction}} (predicting the risk of coronary artery disease); and (9) Myocardial~\cite{misc_myocardial_infarction_complications_579} (classifying individuals for chronic heart failure risks).

We further integrate newer and synthetic datasets that are relatively underexplored and absent from the LLMs’ pre-training corpora: (10) Cultivars, designed to predict the potential grain yield for a specific soybean cultivar\footnote{\texttt{archive.ics.uci.edu/dataset/913}}; (11) NHANES, focused on inferring a person's age group from their data record\footnote{\texttt{archive.ics.uci.edu/dataset/887}}; (12) Sequence-type, for classifying a given number sequence as arithmetic, geometric, Fibonacci, or Collatz; and (13) Solution-mix, tasked with predicting whether the concentration of a mixture from four solutions surpasses a 0.5 threshold. The Cultivars and NHANES datasets were only added to the UCI Machine Learning Repository after September 2023, thereby precluding their inclusion in the pre-training set of the LLM API (gpt-3.5-turbo-0613) used in this work. The last two—Sequence-type and Solution-mix—are synthetic in nature, eliminating the possibility of their memorization by the LLM API and supporting their validity for assessment purposes.

These datasets encompass a broad range in both scale and intricacy, as summarized in Table~\ref{Tab:basic-desc}. Each contains clearly defined names and attribute descriptions. Datasets such as `Communities' and `Myocardial' have attribute counts exceeding 100, which presents challenges, particularly given the limited context length available in current LLMs.

\vspace*{-0.12in}
\paragraph{Baselines.}
For comparative analysis, we benchmark against nine baseline methods. The initial group of three consists of standard approaches for tabular data, namely: (1) Logistic Regression (LogReg), (2) XGBoost~\cite{chen2016xgboost}, and (3) RandomForest~\cite{ho1995random}. The subsequent pair includes methods tailored for scenarios with unlabeled data: (4) SCARF~\cite{bahri2022scarf}, and (5) STUNT~\cite{nam2023stunt}. It should be noted that acquiring extensive unlabeled datasets may pose practical challenges. A contemporary technique, (6) TabPFN~\cite{hollmann2022tabpfn}, utilizes diverse synthetic data distributions for model pretraining. The final three approaches employ LLMs: (7) In-context learning (In-context)~\cite{wei2022emergent}, (8) TABLET~\cite{slack2023tablet}, and (9) TabLLM~\cite{hegselmann2023tabllm}. In-context learning adds few-shot training exemplars directly to the prompt input without model adjustment. TABLET augments the prompt with supplementary cues from an external classifier, using rule-sets and prototypes to bolster inference performance. TabLLM leverages a parameter-efficient adaptation method such as IA3~\cite{liu2022few}, enabling LLM finetuning using only a small set of training examples.

\vspace*{-0.12in}
\paragraph{Implementation details.}
\textsf{FeatLLM} employs GPT-3.5 as the underlying LLM, although its architecture remains flexible to support other LLMs as well (additional results with alternative backbones can be found in Appendix~\ref{sec:backbones}). For LLM inference, we use a temperature of 0.5, and set the top-p parameter to 1, which is the API default. The numbers of ensembles and extraction rules are fixed at 20 and 10, respectively. For insights into the influence of hyper-parameter selections, see Figure~\ref{fig:hyperparameter2} and Appendix~\ref{sec:hyper-parameter}. 

Model training is carried out with the Adam optimizer at a learning rate of 0.01, and proceeds for 200 epochs when fitting the linear model. Epoch selection is guided by $k$-fold cross-validation. In reimplementing baseline methods, in-context learning utilizes GPT-3.5, while TabLLM leverages T0~\cite{sanh2022multitask} as originally prescribed. Note that TabLLM usage is limited to LLMs with openly released checkpoints, and therefore excludes GPT-3.5. Conventional machine learning models, including LogReg, XGBoost, and RandomForest, are tuned using Grid-search in combination with $k$-fold cross-validation. The $k$ parameter is chosen as either 2 or 4, with the constraint that each class must be represented at least once in the training folds. Comprehensive implementation information can be found in Appendix~\ref{sec:baselines}.

\vspace*{-0.12in}
\paragraph{Results.}
Tables~\ref{tab:main1} and \ref{tab:main2} provide a comparative analysis of model performances across various datasets. Each experimental setup is executed three times with varying seeds for selecting training and test splits, and we report both the mean and standard deviation of the resulting AUC scores. When evaluated alongside baseline models, \textsf{FeatLLM} regularly attains either the best or second-best results. Remarkably, our method achieves performance similar to classical tabular models using only 4 shots, while the latter require 32 shots (refer to Fig.~\ref{fig:compares}). Although STUNT and TabPFN yield notable improvements for specific datasets, these gains do not translate to consistent, broad superiority over traditional machine learning approaches. In contrast, large language model-based techniques such as in-context learning, TABLET, or TabLLM are unable to effectively process tabular data with high-dimensional features. Employing both bagging and ensemble mechanisms, our framework demonstrates robust efficacy even with datasets containing a large number of features. Importantly, while the bagging and ensembling strategy we propose can, in principle, be integrated into other LLM-based frameworks, such inclusion would entail a twofold increase in inference computation per sample, greatly intensifying resource demands and reducing practicability.

\subsection{Analysis \& Discussion}
\label{sec:analysis}
\paragraph{Ablation study.} To discern the influence of key architectural choices, we conduct ablation studies targeting the following: (1) \textsf{-Tuning}: removing the process of weight adjustment in the linear model and instead summing binary features to calculate the logit; (2) \textsf{-Ensemble}: eliminating the ensemble stage; (3) \textsf{-Description}: excluding feature description information; and (4) \textsf{-Reasoning}: skipping Step-1 in the reasoning instructions during rule generation.

A summary of AUC score changes from these ablations, averaged across datasets, is shown in Table~\ref{tab:ablation}. For every ablated configuration, a drop in performance is observed, indicating the importance of each component. Notably, the advantage of employing weight tuning is amplified as the number of shots rises, suggesting that tuning is more beneficial when larger data quantities enable better estimation of rule significance. Conversely, feature description and reasoning instruction components exert a stronger influence with fewer shots, highlighting the value of leveraging LLMs’ prior knowledge when labeled data is scarce. Finally, the proposed ensemble approach robustly enhances performance under all evaluated conditions.

\vspace*{-0.12in}
\paragraph{Training \& inference time comparison.} We evaluate the training and inference durations across various approaches. The benchmarking utilizes the Adult dataset, where the training portion comprises 16 shots. All experiments employ a single A100 GPU by default, apart from TabLLM, which leverages four A100 GPUs due to model parallelization requirements. Detailed runtime figures are reported in Table~\ref{tab:cost}. Importantly, \textsf{FeatLLM} achieves notably low inference times, on par with those of traditional techniques. Conversely, other approaches leveraging LLMs—such as In-context learning, TABLET, and TabLLM—entail considerably greater inference durations. When considering training time, \textsf{FeatLLM}'s efficiency is comparable to other LLM-centric methods: it makes just 30 API calls, which is budget-friendly and amenable to parallelization.

\vspace*{-0.12in}
\paragraph{Quality of features generated by \textsf{FeatLLM}.}
We design a straightforward experiment to assess the practical utility of the rule-based features produced by \textsf{FeatLLM}. Specifically, we evaluate the mean AUROC between these constructed features and the target outcomes, subsequently determining the fraction of features for which the AUROC surpasses 0.5 within each dataset. Table~\ref{tab:quality} compiles the experimental findings. The data reveals that most extracted rules are closely aligned with the target variable and yield informative signals for the relevant prediction task (see the column ``Before tuning”). Additionally, during linear model adjustment, those rules deemed uninformative—corresponding to weights below a preset threshold—are systematically excluded from consideration (see ``After tuning”). Here, the threshold was selected to be 0.1, informed by the overall distribution of learned weights.

\vspace*{-0.12in}
\paragraph{Effect of adding spurious correlations.} To evaluate how well different approaches handle noisy, spurious correlations under a low-shot learning scenario, we carried out an experimental study. In this experiment, the Heart dataset was augmented by appending random, irrelevant columns sourced from the Adult dataset—columns that bear no relation to the underlying prediction task. Model performance was monitored as the number of extraneous columns was incrementally increased. For consistency and fairness, all methods were allocated an identical set of 8 labeled samples, without access to any additional unlabeled data; thus, approaches like SCARF and STUNT were not included in this analysis.

Figure~\ref{fig:spurious} depicts how the introduction of spurious correlations affects model accuracy. Among the evaluated methods, \textsf{FeatLLM} exhibited the most stable performance, experiencing minimal decline due to noisy features. This outcome can be attributed to the framework’s reliance on prior knowledge for aligning features with the prediction task during rule extraction, which restricts the generation of rules over unrelated columns and contributes to greater robustness. Empirically, less than 2\% of all extracted rules made use of spurious columns. However, it is noteworthy that when columns from datasets with stronger semantic or causal overlap are appended at random, \textsf{FeatLLM} is susceptible to picking up these unintended correlations and may misconstrue the underlying causal relationships between features and the target variable (refer to Appendix~\ref{sec:spurious-appendix}).

\vspace*{-0.12in}
\paragraph{Hyper-parameter analysis}
The \textsf{FeatLLM} framework primarily depends on two hyper-parameters: the ensemble size, and the quantity of rules per class extracted in each LLM invocation. We performed a series of experiments to evaluate how variations in these parameters influence model effectiveness. The findings indicate that increasing the number of ensembles leads to improved results, although this benefit plateaus after a threshold (approximately 20 ensembles), as illustrated in Fig.~\ref{fig:appendix-hyperparameter1} in the Appendix. Regarding the rules per class, although no significant statistical difference was observed, setting this parameter to 10 provided an optimal trade-off between prompt complexity and predictive performance (refer to Fig.~\ref{fig:hyperparameter2}). We hypothesize that including more rules increases the feature space, potentially boosting information content but also making the model prone to overfitting, especially in low-shot learning contexts.

\section{Conclusion}
In this paper, we introduce \textsf{FeatLLM}, which establishes a new benchmark for few-shot learning on tabular datasets using LLMs. Unlike conventional methods that employ LLMs solely for per-sample predictions, \textsf{FeatLLM} leverages LLMs as feature engineering tools by synthesizing rule-based criteria. With its combination of rule induction and an ensemble strategy based on bagging, \textsf{FeatLLM} offers a solution that keeps inference costs low, operates using only API calls—no training required—and addresses the limitations imposed by data dimensionality. Our results demonstrate that \textsf{FeatLLM} achieves robust predictive outcomes, offering an effective and efficient method for deploying LLMs on practical tabular learning problems.

The scope of \textsf{FeatLLM} is currently focused on scenarios with extremely limited labeled data. Moving forward, we aim to extend its methodological arsenal to enable feature creation for datasets containing more examples, experiment with diverse feature representations beyond rules, and strengthen interpretability, thereby facilitating broader applicability to additional machine learning problems.

\section{Broader Impact}
The framework we introduce, \textsf{FeatLLM}, is developed to harness the prior knowledge encapsulated in LLMs, pushing beyond what conventional supervised tabular learning can achieve—especially under data-scarce conditions. By leveraging pretrained knowledge, our work addresses overfitting challenges arising from small sample sizes and provides an approach to learning that is significantly more data-efficient. \textsf{FeatLLM} is particularly advantageous for industries like finance and healthcare, where collecting labeled examples can be prohibitively costly or infeasible, and LLMs’ domain-specific prior knowledge—acquired through extensive text pretraining—has added value. However, one of the essential directions for future exploration involves examining the implications of societal bias and misinformation that may reside in LLMs’ knowledge bases, as these biases can substantially impact model predictions, at times outweighing the influence of observed labeled data.

\end{document}